\begin{abstract}
지금의 VLA는 3--8B VLM 하나가 계획·판단·명령·움직임을 모두 떠안는다. 그 크기의 모델에게는 무거운 짐이고, VLM에서 행동 전문가로 넘어가며 정보도 잃는다. Fly-by-Wire는 이 짐을 나눈다. 학습 가능한 크기의 상위 VLM(Qwen3.5-35B-A3B, LoRA)이 계획·판단·명령을 모두 맡고, 머리 영상 위의 점과 높이 의도로 로봇을 조종한다. 목표 좌표는 모델이 아니라 코드가 센서 깊이로 계산한다. 상위는 결합을 위해 다시 학습하지 않는다: 규칙 코드 어댑터(학습 없음)가 상위 출력(점·높이 의도·그리퍼 의도·평가)을 목적지 xyz, 대상 물체 마스크, 단계, 그리퍼 허용, 단계별 끌림 세기(운반 중 약, 잡기·놓기 근처 강 — 딱딱한 범위 대신 목적지로 당김, 넓은 안전 범위만 자름)로 바꾼다. 하위 JCR(Joystick-Conditioned Reflex, 조이스틱 반사 정책: 언어·계획 없이 그 끌림을 따르고, 잡기·놓기 시점과 접촉만 판단하는 반사 정책)는 이 끌림을 좇는 좁은 조이스틱이며, 이상 신호는 어댑터가 문장으로 바꿔 상위의 다음 호출에 실어 보낸다. 두 모듈이 각자 최대 성능을 내도 연결에서 성능이 샐 수 있으므로, 우리는 연결 손실을 같은 시드의 2$\times$2 참값 교체(어댑터 포함)로 정의하고 5\,\%p 이하를 합격 기준으로 둔다. 재질의는 도착 뒤(상위가 학습한 장면 그대로)이며, 상위가 1--3\,s 생각하는 동안 JCR가 정렬·잡기·들기 시작을 이어가 로봇이 멈추지 않는다. 상위는 동결하고, JCR만 어댑터가 실제로 내는 명령(상위 흉내 명령도 어댑터 통과) 위에서 학습해 맞물린다. 지금까지 사전 등록 실험에서 깊이 점 출력은 본 적 없는 탁자 높이에서 \tentative{3.5\,mm}를 유지했고(거리 숫자 출력은 \tentative{31--71\,mm}), 공개 로봇 데이터의 점 라벨은 새 물체 오차를 \tentative{85.3}에서 \tentative{4.0\,mm}로 줄였다. 시뮬 참값 명령을 스크립트 실행기가 따르면 \tentative{40/40}이었지만 기존 VLA를 조이스틱으로 쓰면 \tentative{1/40}이었다. 폐루프 상태 분포와 근접 미세 조종을 배우지 않은 탓이며, 그래서 JCR는 자기 롤아웃에 시뮬 참값을 다시 붙여 반복 학습한다.
\end{abstract}
